\section{Open Questions (5)}

\subsection*{Q1. Ion-species specificity of the Akkermansia bloom}
\textbf{Basis.} The deposited SRA SRP098151 contains only $^{16}$O-irradiated mice. The paper is framed as ``space-type radiation'' but does not test cross-ion generalizability. Without cross-ion data, the mechanism cannot be attributed to LET / track-structure vs simply $^{16}$O chemistry, and the ``space radiation'' framing overreaches the empirical evidence. \\
\textbf{Next steps.} Cross-search NASA GeneLab / OSDR and NCBI SRA for 16S amplicon datasets from $^{56}$Fe, $^{28}$Si, or proton whole-body irradiation of C57BL/6 mice at matched 0.1--1~Gy doses. Re-run the identical vsearch + SILVA 138 pipeline (\texttt{scripts/04-07}) on any matching cohort. Compare Akkermansia relative abundance and Bray-Curtis PERMANOVA F-statistics vs the $^{16}$O baseline. If no matching public data exist, this is a concrete experimental proposal for a follow-up NSRL beam campaign.

\subsection*{Q2. Causality vs correlation for the Akkermansia bloom}
\textbf{Basis.} The paper (and our replication) is strictly observational: 16S community composition shifts with dose and time. No germ-free re-derivation, gnotobiotic transfer, or \emph{A.~muciniphila} mono-colonization experiment is performed. The observational-vs-mechanistic distinction is preserved but the causal question is untouched. \\
\textbf{Next steps.} Design a two-arm follow-up: (a) germ-free C57BL/6 mice colonized with pre- vs post-0.1-Gy-$^{16}$O donor cecal contents, phenotyped for tight-junction integrity (ZO-1/occludin IHC), FITC-dextran gut permeability, and serum LPS/LBP; (b) SPF mice given oral \emph{A.~muciniphila} (BAA-835) supplementation before 0.1~Gy $^{16}$O and compared to vehicle for the same readouts. Free-endpoint tool support: search NCBI GEO/ArrayExpress for existing germ-free + heavy-ion datasets to see if partial evidence already exists.

\subsection*{Q3. Mouse-to-human translation}
\textbf{Basis.} Murine gut microbiome (dominated by Muribaculaceae/S24-7, high Akkermansia baseline capacity) differs qualitatively from human (Bacteroides/Prevotella/Firmicutes-dominant). Astronaut microbiome studies (JAXA Kibo, NASA Twins Study) show shifts but with different taxonomic signatures. Direct translation of the 0.1~Gy Akkermansia bloom to humans is unsupported. \\
\textbf{Next steps.} Re-analyse public astronaut 16S/metagenomic data: NASA Twins Study (Garrett-Bakelman 2019, \emph{Science}) fecal 16S, JAXA Kibo pre/post-flight cohorts (if available via DDBJ), and ISS crewmember stool time-series. Apply the same alpha/beta/differential-abundance pipeline as this replication. Ask: is there any dose-response or flight-duration effect on Akkermansia or Verrucomicrobiota fraction? Cross-reference to GCR dose reconstruction from active dosimetry (Zeitlin et al.).

\subsection*{Q4. Cage-effect confounding}
\textbf{Basis.} Mice are coprophagic and share microbiota within a cage; dose group and cage are potentially confounded if the 80 mice were housed in fewer than 8 cages with dose-per-cage segregation. The paper does not describe cage assignment or use mixed-effect models with cage as a random effect. Our replication inherits this ambiguity. \\
\textbf{Next steps.} Contact corresponding authors (Cheema AK, Braun J) to request per-mouse cage IDs; if obtained, re-fit PERMANOVA with cage as a stratum and DESeq2 with cage as a random effect (\texttt{variancePartition} or \texttt{glmmTMB}). Estimate ICC(cage). If cage explains $>20\%$ of variance orthogonal to dose, headline effects need reweighting. As a partial substitute, re-analyze at the level of the 8 dose$\times$time cells and treat within-cell replicate correlation empirically via bootstrap.

\subsection*{Q5. Prebiotic / A.~muciniphila countermeasure viability}
\textbf{Basis.} The paper documents the perturbation but proposes no countermeasure. Akkermansia is a well-studied mucin-degrader linked to gut barrier and metabolic health; prebiotics that promote it are plausible protective agents. This is directly actionable for NASA crew health. \\
\textbf{Next steps.} Literature triangulation: (i) systematic search of PubMed for ``(inulin OR FOS OR prebiotic OR pasteurized Akkermansia) AND (radiation OR heavy ion) AND (mouse OR gut)''; (ii) query NASA GeneLab for any dietary-intervention arm in existing radiation-microbiome experiments; (iii) if none exist, this is a straightforward pre-registered proposal for a 4-arm mouse study (vehicle vs inulin vs live A.m.\ vs pasteurized A.m., all $\times$ 0.1~Gy $^{16}$O vs sham) with 30-day 16S + gut-barrier + LC-MS metabolome as endpoints, deliberately including the metabolomics arm this paper failed to deposit.
